% Source: https://artofproblemsolving.com/community/c3h1781571_help_needed_counting_problem
% Stable identity: aops-topic-1781571-q01
The digits are $1,2,2,3,3,4,5$: the digits $2$ and $3$ occur twice and $1,4,5$ once. There is no $0$, so every arrangement is a genuine number.

(a) The number of arrangements of a multiset is $\dfrac{7!}{2!\,2!}=\dfrac{5040}{4}=\boxed{1260}$.

(b) Classify the five chosen digits by how many repeated pairs they contain.
\begin{itemize}
\item Five different digits: they must be $1,2,3,4,5$, giving $5!=120$ numbers.
\item Exactly one pair ($22$ or $33$) and three other different digits chosen from the remaining four values: $2\cdot\dbinom43=8$ choices of digits, each with $\dfrac{5!}{2!}=60$ arrangements, giving $480$.
\item Both pairs $22$, $33$ and one of $1,4,5$: $3$ choices, each with $\dfrac{5!}{2!\,2!}=30$ arrangements, giving $90$.
\end{itemize}
The total is $120+480+90=\boxed{690}$.
